\chapter{CHAPTER 7: CONCLUSION AND FUTURE
WORK}\label{chapter-7-conclusion-and-future-work}

\section{7.1 Chapter Overview}\label{chapter-overview}

This final chapter synthesizes the outcomes of the SentinelAQ research
project. Section 7.2 evaluates the completion of the established
research objectives. Section 7.3 provides the overall conclusion of the
study, directly answering the core research questions. Finally, Section
7.4 outlines strategic, forward-looking future work, identifying
specific avenues for expanding this research beyond its current scope.

\section{7.2 Research Summary and Objective
Evaluation}\label{research-summary-and-objective-evaluation}

This research set out to design, develop, and evaluate a multi-horizon
machine learning forecasting framework tailored for the distinct urban
topographies of Sri Lanka. A review of the methodology and results
confirms that all seven SMART objectives established in Chapter 1 were
successfully executed: 1. \textbf{Data Synchronization:} A multi-modal
dataset was successfully constructed by fusing PurpleAir telemetry with
Open-Meteo meteorological proxies. 2. \textbf{Feature Engineering:} The
critical 48-dimensional feature vector, utilizing rolling means and
cyclical encodings, was engineered to capture non-linear accumulation.
3. \textbf{Model Training:} SARIMAX, XGBoost, and LSTM models were
rigorously trained. 4. \textbf{Multi-Horizon Evaluation:} The models
were mathematically benchmarked across 1h to 48h horizons, proving
topographical variances between Colombo and Kandy. 5.
\textbf{Explainability:} SHAP was integrated, successfully
deconstructing the predictive ``black box'' into actionable atmospheric
insights. 6. \textbf{Pipeline Automation:} The inference logic was
wrapped into an automated Python backend synchronized with Google
Firebase. 7. \textbf{Mobile Deployment:} The React Native SentinelAQ
application was successfully deployed, visualizing real-time forecasts
and triggering threshold-based push notifications.

\section{7.3 Overall Conclusion}\label{overall-conclusion}

The primary problem addressed by this study was the lack of reliable,
topography-aware, and multi-horizon PM2.5 forecasting in Sri Lanka. By
iteratively developing the SentinelAQ architecture, this research
yielded three definitive conclusions that contribute to the domain of
environmental computer science:

First, the research concluded that reliance on optical/chemical
satellite data (e.g., Sentinel-5P) is currently unviable for
high-frequency, short-term forecasting in tropical microclimates. Severe
cloud occlusion over mountainous regions like Kandy necessitates a
structural reliance on robust ground-level meteorological proxies to
ensure data continuity for machine learning pipelines.

Second, the algorithmic benchmarking proved that tree-based ensemble
methods (XGBoost) significantly outperform sequential deep learning
(LSTM) for multi-horizon environmental forecasting on sparse datasets.
By engineering time-lags and rolling statistical means into a tabular
format, XGBoost maintained predictive stability up to 48 hours in the
future, whereas statistical models (SARIMAX) collapsed after 1 hour, and
LSTMs systematically overfit. Furthermore, XGBoost proved vastly
superior in implementation speed and computational efficiency, making it
the optimal choice for cloud-based deployment.

Finally, the integration of Explainable AI (SHAP) mathematically proved
that urban topography dictates pollution dynamics. The system revealed
that coastal pollution in Colombo is primarily driven by macro-seasonal
wind patterns (transboundary dispersion), whereas valley pollution in
Kandy is driven by short-term localized accumulation and thermal
inversions (valley trapping). Ultimately, SentinelAQ demonstrates that
proactive, highly accurate, and transparent public health alerting is
computationally feasible in developing nations without requiring massive
state-funded sensor infrastructure.

\section{7.4 Strategic Future Work}\label{strategic-future-work}

The development of SentinelAQ establishes a robust foundational
architecture. To advance this domain further, future research must build
\emph{upon} this infrastructure to address broader environmental
challenges.

\subsection{7.4.1 Spatial Expansion via Transfer Learning
(Lesser-Covered
Districts)}\label{spatial-expansion-via-transfer-learning-lesser-covered-districts}

The current research scope deliberately targeted the two primary
economic hubs (Colombo and Kandy) where dense historical sensor data was
available. A critical avenue for future work involves deploying the
SentinelAQ architecture to rapidly industrializing, lesser-covered
districts (such as the Hambantota Port Zone or Jaffna) where historical
data is nonexistent. Future research should investigate the application
of \textbf{Transfer Learning}. By taking the pre-trained XGBoost weights
from the Colombo/Kandy models and applying them to new geographical
nodes utilizing only real-time weather API inputs, researchers could
effectively deploy ``virtual sensor nodes.'' This would allow the system
to rapidly scale and map the entire island's air quality without the
prohibitive cost of installing physical hardware in every district.

\subsection{7.4.2 Holistic AQI Scale Integration (Uncovered Chemical
Elements)}\label{holistic-aqi-scale-integration-uncovered-chemical-elements}

To address the most immediate threat to public health, this research
specifically isolated and modeled PM2.5, which is the primary driver of
Sri Lanka's pollution spikes. However, the true US EPA Air Quality Index
(AQI) is a composite scale comprising multiple criteria pollutants. As
Sri Lanka's industrial manufacturing and energy sectors expand, future
iterations of the SentinelAQ inference pipeline must be upgraded to
dynamically ingest and forecast uncovered, secondary chemical
elements---specifically \textbf{Surface Ozone (O3)} and \textbf{Sulfur
Dioxide (SO2)}. Integrating these distinct chemical vectors into the
existing XGBoost feature matrix will transition the architecture from a
targeted PM2.5 alerting tool into a fully holistic, EPA-standardized AQI
forecasting engine, providing a more comprehensive defense for public
health.
